\documentclass[10pt,twocolumn,letterpaper]{article}

\usepackage{times}
\usepackage{epsfig}
\usepackage{graphicx}
\usepackage{amsmath}
\usepackage{amssymb}
\usepackage{booktabs}
\usepackage{hyperref}

\title{Structural Invariance and Domain Shift Resilience in Phishing Detection:\\A Comparative Analysis of Deep Learning, Bagging, and Incremental Boosting}

\author{
Anonymous Authors\\
\tt\small research@phishing-framework.org
}

\begin{document}

\maketitle

\begin{abstract}
As phishing templates evolve to evade detection, traditional machine learning models suffer catastrophic accuracy degradation due to domain shift. In this paper, we evaluate the resilience of various model architectures—including Convolutional Neural Networks (CNN), Random Forests (RF), and eXtreme Gradient Boosting (XGBoost)—across a 5-year longitudinal dataset spanning 45,000 historical records and 40,000 modern "zero-day" records. We demonstrate that while Deep Learning NLP models achieve near-perfect (99.91\%) in-domain accuracy, they collapse to $\sim$49\% during domain shift due to vocabulary misalignment. To solve this, we engineer a strictly mathematical Structural Feature space (DOM topology, CSS hacks, lexical ratios). Furthermore, we prove that static structural models still suffer from covariate shift. Ultimately, we propose a Continuous Machine Learning pipeline utilizing XGBoost, which successfully self-heals from an 11.12\% zero-day failure rate to 99.74\% accuracy by incrementally learning as few as 200 target-domain samples, massively outperforming standard Data Mixing retraining methodologies.
\end{abstract}

\section{Introduction}
Phishing detection is an adversarial cat-and-mouse game. Security researchers train models on historical phishing templates, but scammers continuously mutate their DOM structures, vocabulary, and hosting infrastructure to bypass static filters. This phenomenon, known as Domain Shift or Concept Drift, is the primary reason why academic models with 99\% accuracy fail in live industrial deployment.

This paper systematically investigates the failure points of static models when exposed to multi-year domain shifts. We analyze the "Curse of Dimensionality" in structural feature engineering, the Data Dilution effect in Few-Shot retraining, and the mathematical boundaries of Incremental Learning. 

\section{Methodology}

\subsection{Datasets}
To strictly measure domain shift, we utilized three chronologically distinct datasets:
\begin{enumerate}
    \item \textbf{Historical Base (2021):} 45,373 samples containing generic historical HTML and URL strings.
    \item \textbf{PhreshPhish Dataset (Modern):} 40,000 heavily imbalanced, sophisticated phishing templates representing modern evasion tactics.
    \item \textbf{Zero-Day Stream (2026):} 5,841 freshly crawled out-of-distribution (OOD) domains to test terminal real-world resilience.
\end{enumerate}

\subsection{Structural Invariance Engineering}
Rather than relying on NLP tokens (which change rapidly), we engineered a \textit{StructuralProcessor} to extract 44 invariant mathematical properties. These include DOM topology (tag sequence N-grams), CSS visibility anomalies (e.g., \texttt{display:none}), and lexical ratios (e.g., URL entropy, digit density). Recursive Feature Elimination (RFE) was used to prune noise and isolate the most stable invariants.

\section{Experimental Results}

\subsection{The Deep Learning Collapse}
Our initial baseline evaluated the WebPhish CNN. Evaluated In-Domain, it achieved 99.91\% accuracy. However, when evaluated on the 2026 Zero-Day dataset, accuracy collapsed to sub-50\%. This proved that NLP-based Deep Learning networks merely memorize temporal vocabulary strings (e.g., "paypal\_login") and possess zero structural resilience to domain shift.

\subsection{Static Structural Models (RF vs XGBoost)}
Transitioning to structural invariants, we trained Random Forest and XGBoost architectures. 
\begin{itemize}
    \item \textbf{XGBoost (In-Domain):} 97.25\% Accuracy.
    \item \textbf{Random Forest (In-Domain):} 96.53\% Accuracy.
\end{itemize}
However, when exposed to 5-year domain drift, static thresholds aged poorly, resulting in catastrophic failure rates (11.12\% accuracy on pristine Zero-Day data) due to Covariate Shift.

\subsection{The Incremental Learning Breakthrough}
To fix Covariate Shift without destroying the base model, we implemented Incremental Boosting via XGBoost's \texttt{warm\_start} tree appending. 
When streaming the 2026 Zero-Day data into the frozen model:
\begin{itemize}
    \item \textbf{0 Samples:} 11.12\% (Collapse)
    \item \textbf{200 Samples:} 99.74\% (Instant Recovery)
    \item \textbf{800 Samples:} 100.00\% (Perfect Calibration)
\end{itemize}
XGBoost required a mere 200 samples to recalibrate its internal split thresholds to the new domain.

\subsection{The Limits of Incremental Adaptation}
We hypothesized whether Incremental Learning could completely replace full retraining if the target dataset is structurally deeper than the base dataset. Streaming 32,000 PhreshPhish samples into the simple 11k base model plateaued permanently at $\sim$90.30\%. This proved that Incremental Learning handles threshold calibration but cannot fundamentally invent deep architectural splits.

\subsection{The Data Dilution Effect in Retraining}
When forced to fully retrain, we injected target-domain samples into the 45k historical dataset using Random Forest. Injecting 5,000 target samples only yielded 92.19\% accuracy. The massive historical data mathematically "diluted" the new samples during bagging, proving that raw data mixing is inefficient for domain adaptation.

\section{Conclusion}
Our research confirms that XGBoost is the ultimate operational architecture for Phishing Detection. While Random Forest is 1.5x faster to retrain from scratch, XGBoost possesses the unique capability to execute Continuous Machine Learning. By streaming extremely small batches of zero-day data into XGBoost, security systems can exponentially adapt to domain shifts in milliseconds, entirely bypassing the Data Dilution effect and the immense computational costs of Deep Learning retraining.

\end{document}
